\RequirePackage{plautopatch}
\documentclass[platex,dvipdfmx,a4paper]{jsarticle}

\usepackage{amsmath}
\usepackage{amsfonts}
\usepackage{amssymb}
\usepackage[dvipdfmx]{graphicx}
\usepackage{float}
\usepackage{tikz}
\usepackage{caption}
\usepackage[%
 dvipdfmx,%
 setpagesize=false,%
 bookmarks=true,%
 bookmarksdepth=tocdepth,%
 bookmarksnumbered=true,%
 colorlinks=false,%
 pdftitle={物理数学II期末試験解答},%
 pdfsubject={2023年度から2025年度},%
 pdfauthor={},%
 pdfkeywords={}%
]{hyperref}
\usepackage{pxjahyper}

\newcommand{\hb}[2]{\frac{\partial #1}{\partial #2}}
\newcommand{\bbn}[2]{\frac{d #1}{d #2}}
\newcommand{\nhb}[3]{\frac{\partial^{#1} #2}{\partial #3^{#1}}}
\newcommand{\nbbn}[3]{\frac{d^{#1} #2}{d #3^{#1}}}
\newcommand{\vectxy}[2]{\begin{pmatrix}#1\\#2\end{pmatrix}}
\newcommand{\vectxyz}[3]{\begin{pmatrix}#1\\#2\\#3\end{pmatrix}}
\newcommand{\vect}[1]{\boldsymbol{\mathrm{#1}}}
\newcommand{\ssum}[3]{\sum_{#1=#2}^{#3}}

\begin{document}

\section*{物理数学II期末試験解答}

\section*{2023年度}

\subsection*{問題1}

区間 $-\pi\leq x<\pi$ における $f(x)=x$ の Fourier 級数を求める。周期 $2\pi$ の Fourier 級数を：
\begin{align*}
  f(x)
  &=\frac{a_0}{2}+\sum_{n=1}^{\infty}\left(a_n\cos nx+b_n\sin nx\right)
\end{align*}
とする。$f(x)=x$ は奇関数であるため、$a_0=a_n=0$ である。正弦係数は：
\begin{align*}
  b_n
  &=\frac{1}{\pi}\int_{-\pi}^{\pi}x\sin nx\,dx
   =\frac{2}{\pi}\int_0^\pi x\sin nx\,dx。
\end{align*}
部分積分を用いる：
\begin{align*}
  \int_0^\pi x\sin nx\,dx
  &=\left[-\frac{x\cos nx}{n}\right]_0^\pi
    +\frac{1}{n}\int_0^\pi\cos nx\,dx \\
  &=-\frac{\pi(-1)^n}{n}。
\end{align*}
したがって：
\begin{align*}
  b_n
  &=\frac{2(-1)^{n+1}}{n}。
\end{align*}
以上より：
\begin{align*}
  x
  &=2\sum_{n=1}^{\infty}\frac{(-1)^{n+1}}{n}\sin nx
  \qquad (-\pi<x<\pi)。
\end{align*}

\subsection*{問題2}

領域 $-\infty<x<\infty$, $0\leq t<\infty$ において、熱方程式：
\begin{align*}
  \hb{u(x,t)}{t}
  &=\kappa\nhb{2}{u(x,t)}{x}
\end{align*}
を、初期条件 $u(x,0)=\delta(x)$ の下で解く。

\subsubsection*{1}

問題文に従い、逆 Fourier 変換を：
\begin{align*}
  u(x,t)
  &=\frac{1}{2\pi}\int_{-\infty}^{\infty}U(k,t)e^{ikx}\,dk
\end{align*}
とする。これを熱方程式へ代入する：
\begin{align*}
  \frac{1}{2\pi}\int_{-\infty}^{\infty}\hb{U(k,t)}{t}e^{ikx}\,dk
  &=-\frac{\kappa}{2\pi}\int_{-\infty}^{\infty}k^2U(k,t)e^{ikx}\,dk。
\end{align*}
Fourier 変換の一意性から：
\begin{align*}
  \hb{U(k,t)}{t}
  &=-\kappa k^2U(k,t)。
\end{align*}
また、$t=0$ において：
\begin{align*}
  \delta(x)
  &=\frac{1}{2\pi}\int_{-\infty}^{\infty}U(k,0)e^{ikx}\,dk
\end{align*}
である。デルタ関数の Fourier 表示と比較すると：
\begin{align*}
  U(k,0)&=1。
\end{align*}

\subsubsection*{2}

得られた微分方程式を変数分離する：
\begin{align*}
  \frac{1}{U}\hb{U}{t}
  &=-\kappa k^2。
\end{align*}
$t$ について積分し、初期条件を課す：
\begin{align*}
  U(k,t)
  &=A(k)e^{-\kappa k^2t}
   =e^{-\kappa k^2t}。
\end{align*}

\subsubsection*{3}

逆 Fourier 変換へ戻す：
\begin{align*}
  u(x,t)
  &=\frac{1}{2\pi}\int_{-\infty}^{\infty}e^{-\kappa k^2t+ikx}\,dk。
\end{align*}
指数部分を平方完成する：
\begin{align*}
  -\kappa tk^2+ikx
  &=-\kappa t\left(k-\frac{ix}{2\kappa t}\right)^2-\frac{x^2}{4\kappa t}。
\end{align*}
積分路を平行移動した後、Gauss 積分を用いる：
\begin{align*}
  u(x,t)
  &=\frac{e^{-x^2/(4\kappa t)}}{2\pi}
    \int_{-\infty}^{\infty}e^{-\kappa tk^2}\,dk \\
  &=\frac{1}{\sqrt{4\pi\kappa t}}
    \exp\left(-\frac{x^2}{4\kappa t}\right)。
\end{align*}
以上より、求める解は：
\begin{align*}
  u(x,t)
  &=\frac{1}{\sqrt{4\pi\kappa t}}
    \exp\left(-\frac{x^2}{4\kappa t}\right)
  \qquad (t>0)
\end{align*}
である。

\subsection*{問題3}

Rodrigues の公式：
\begin{align*}
  P_n(x)
  &=\frac{1}{2^nn!}\nbbn{n}{}{x}(x^2-1)^n
\end{align*}
を用いる。低次の Legendre 多項式は：
\begin{align*}
  P_0(x)&=1,\\
  P_1(x)&=\frac{1}{2}\bbn{}{x}(x^2-1)=x,\\
  P_2(x)&=\frac{1}{8}\nbbn{2}{}{x}(x^2-1)^2=\frac{1}{2}(3x^2-1),\\
  P_3(x)&=\frac{1}{48}\nbbn{3}{}{x}(x^2-1)^3=\frac{1}{2}(5x^3-3x)。
\end{align*}
次に、$x^3$ は奇関数であるため：
\begin{align*}
  x^3
  &=aP_1(x)+bP_3(x)
\end{align*}
と置ける。具体形を代入して係数を比較する：
\begin{align*}
  x^3
  &=ax+\frac{b}{2}(5x^3-3x)。
\end{align*}
したがって：
\begin{align*}
  \frac{5b}{2}&=1,
  &a-\frac{3b}{2}&=0
\end{align*}
であるから、$b=2/5$, $a=3/5$ を得る。以上より：
\begin{align*}
  x^3
  &=\frac{3}{5}P_1(x)+\frac{2}{5}P_3(x)。
\end{align*}

\subsection*{問題4}

Bessel 関数を：
\begin{align*}
  J_n(z)
  &=\frac{1}{\pi}\int_0^\pi\cos(z\sin\theta-n\theta)\,d\theta,
  \qquad n=0,\pm1,\pm2,\ldots
\end{align*}
で定義する。

\subsubsection*{1}

$n$ を $-n$ に置き換え、$\theta=\pi-\varphi$ と変数変換する：
\begin{align*}
  J_{-n}(z)
  &=\frac{1}{\pi}\int_0^\pi\cos(z\sin\varphi+n\pi-n\varphi)\,d\varphi \\
  &=(-1)^nJ_n(z)。
\end{align*}

\subsubsection*{2}

定義式を $z$ で微分する：
\begin{align*}
  \bbn{J_n(z)}{z}
  &=-\frac{1}{\pi}\int_0^\pi\sin\theta\sin(z\sin\theta-n\theta)\,d\theta。
\end{align*}
一方、三角関数の加法定理から：
\begin{align*}
  J_{n-1}(z)-J_{n+1}(z)
  &=-\frac{2}{\pi}\int_0^\pi\sin\theta\sin(z\sin\theta-n\theta)\,d\theta。
\end{align*}
したがって：
\begin{align*}
  \bbn{J_n(z)}{z}
  &=\frac{1}{2}\left(J_{n-1}(z)-J_{n+1}(z)\right)。
\end{align*}

\subsubsection*{3}

次の全微分を用いる：
\begin{align*}
  \bbn{}{\theta}\sin(z\sin\theta-n\theta)
  &=(z\cos\theta-n)\cos(z\sin\theta-n\theta)。
\end{align*}
両辺を $0$ から $\pi$ まで積分すると境界項は零であるため：
\begin{align*}
  nJ_n(z)
  &=\frac{z}{\pi}\int_0^\pi\cos\theta\cos(z\sin\theta-n\theta)\,d\theta。
\end{align*}
また、加法定理から：
\begin{align*}
  J_{n-1}(z)+J_{n+1}(z)
  &=\frac{2}{\pi}\int_0^\pi\cos\theta\cos(z\sin\theta-n\theta)\,d\theta。
\end{align*}
よって：
\begin{align*}
  \frac{n}{z}J_n(z)
  &=\frac{1}{2}\left(J_{n-1}(z)+J_{n+1}(z)\right)。
\end{align*}

\subsubsection*{4}

問2と問3から：
\begin{align*}
  J_{n-1}(z)&=J_n'(z)+\frac{n}{z}J_n(z),\\
  J_{n+1}(z)&=-J_n'(z)+\frac{n}{z}J_n(z)。
\end{align*}
問2を微分し、隣接次数にも問2を適用する：
\begin{align*}
  J_n''
  &=\frac{1}{2}\left(J_{n-1}'-J_{n+1}'\right) \\
  &=\frac{1}{4}\left(J_{n-2}-2J_n+J_{n+2}\right)。
\end{align*}
問3から：
\begin{align*}
  J_{n-2}+J_n&=\frac{2(n-1)}{z}J_{n-1},\\
  J_n+J_{n+2}&=\frac{2(n+1)}{z}J_{n+1}
\end{align*}
である。これらを代入する：
\begin{align*}
  J_n''
  &=-J_n+\frac{n-1}{2z}J_{n-1}+\frac{n+1}{2z}J_{n+1} \\
  &=-J_n-\frac{1}{z}J_n'+\frac{n^2}{z^2}J_n。
\end{align*}
したがって：
\begin{align*}
  \nbbn{2}{J_n}{z}
  +\frac{1}{z}\bbn{J_n}{z}
  +\left(1-\frac{n^2}{z^2}\right)J_n
  &=0。
\end{align*}

\section*{2024年度}

\subsection*{問題1}

\subsubsection*{1}

$f(x)=|x|$ は偶関数であるため、Fourier 級数には余弦項のみが現れる。定数項は：
\begin{align*}
  a_0
  &=\frac{2}{\pi}\int_0^\pi x\,dx
   =\pi。
\end{align*}
$n\geq1$ に対して：
\begin{align*}
  a_n
  &=\frac{2}{\pi}\int_0^\pi x\cos nx\,dx \\
  &=\frac{2}{\pi n^2}\left((-1)^n-1\right)。
\end{align*}
したがって、偶数 $n$ では $a_n=0$、奇数 $n$ では $a_n=-4/(\pi n^2)$ である。以上より：
\begin{align*}
  |x|
  &=\frac{\pi}{2}-\frac{4}{\pi}\sum_{m=0}^{\infty}
    \frac{\cos(2m+1)x}{(2m+1)^2}。
\end{align*}

\subsubsection*{2}

Fourier 級数を：
\begin{align*}
  f(x)&=\sum_{m=-\infty}^{\infty}a_me^{imx},
  &g(x)&=\sum_{n=-\infty}^{\infty}b_ne^{inx}
\end{align*}
とする。畳み込みの定義へ代入する：
\begin{align*}
  h(x)
  &=\frac{1}{2\pi}\int_{-\pi}^{\pi}f(y)g(x-y)\,dy \\
  &=\sum_{m,n}a_mb_ne^{inx}\frac{1}{2\pi}
    \int_{-\pi}^{\pi}e^{i(m-n)y}\,dy。
\end{align*}
複素指数関数の直交関係を用いる：
\begin{align*}
  \frac{1}{2\pi}\int_{-\pi}^{\pi}e^{i(m-n)y}\,dy
  &=\delta_{mn}。
\end{align*}
したがって：
\begin{align*}
  h(x)
  &=\sum_{n=-\infty}^{\infty}a_nb_ne^{inx}。
\end{align*}
以上より：
\begin{align*}
  c_n&=a_nb_n。
\end{align*}

\subsection*{問題2}

\subsubsection*{1}

任意の滑らかな試験関数 $\varphi(x)$ に対して、部分積分を用いる：
\begin{align*}
  \int_{-\infty}^{\infty}\bbn{\Theta(x-x')}{x}\varphi(x)\,dx
  &=-\int_{-\infty}^{\infty}\Theta(x-x')\varphi'(x)\,dx \\
  &=-\int_{x'}^\infty\varphi'(x)\,dx
   =\varphi(x')。
\end{align*}
これはデルタ関数の定義と一致するため、超関数の意味で：
\begin{align*}
  \bbn{\Theta(x-x')}{x}
  &=\delta(x-x')。
\end{align*}

\subsubsection*{2}

$x<x'$ で零となる形として：
\begin{align*}
  G(x,x')
  &=-(x-x')\Theta(x-x')
\end{align*}
と置く。超関数の公式 $(x-x')\delta(x-x')=0$ を用いる：
\begin{align*}
  \hb{G(x,x')}{x}
  &=-\Theta(x-x'),\\
  \nhb{2}{G(x,x')}{x}
  &=-\delta(x-x')。
\end{align*}
また、$x<x'$ では $\Theta(x-x')=0$ である。したがって：
\begin{align*}
  G(x,x')
  &=-(x-x')\Theta(x-x')
\end{align*}
が求める Green 関数である。

\subsection*{問題3}

生成関数を：
\begin{align*}
  G(x,t)
  &=\frac{1}{\sqrt{1-2tx+t^2}}
   =\sum_{n=0}^{\infty}P_n(x)t^n
\end{align*}
とする。

\subsubsection*{1}

二項展開を $t^2$ まで行う：
\begin{align*}
  G(x,t)
  &=1+xt+\frac{1}{2}(3x^2-1)t^2+O(t^3)。
\end{align*}
よって：
\begin{align*}
  P_0(x)&=1,
  &P_1(x)&=x,
  &P_2(x)&=\frac{1}{2}(3x^2-1)。
\end{align*}

\subsubsection*{2}

生成関数を $t$ で微分する：
\begin{align*}
  \hb{G}{t}
  &=(x-t)(1-2tx+t^2)^{-3/2}。
\end{align*}
したがって：
\begin{align*}
  (1-2tx+t^2)\hb{G}{t}
  &=(x-t)G。
\end{align*}
級数表示を代入し、$t^n$ の係数を比較する：
\begin{align*}
  (n+1)P_{n+1}(x)
  &=(2n+1)xP_n(x)-nP_{n-1}(x)。
\end{align*}

\subsubsection*{3}

生成関数の偏微分を計算する：
\begin{align*}
  \hb{G}{x}
  &=t(1-2tx+t^2)^{-3/2},\\
  \nhb{2}{G}{x}
  &=3t^2(1-2tx+t^2)^{-5/2},\\
  \hb{G}{t}
  &=(x-t)(1-2tx+t^2)^{-3/2},\\
  \nhb{2}{G}{t}
  &=-(1-2tx+t^2)^{-3/2}
    +3(x-t)^2(1-2tx+t^2)^{-5/2}。
\end{align*}
これらを代入し、分母をそろえて整理すると：
\begin{align*}
  (1-x^2)\nhb{2}{G}{x}-2x\hb{G}{x}
  &=-t^2\nhb{2}{G}{t}-2t\hb{G}{t}
\end{align*}
を得る。なお、もし問題用紙の右辺に負号が印刷されていない場合、その式は $P_1(x)=x$ を代入しただけでも成立しない。生成関数から直接得られる恒等式は上式である。

\subsubsection*{4}

級数表示を問3の偏微分方程式へ代入する。左辺は：
\begin{align*}
  (1-x^2)\nhb{2}{G}{x}-2x\hb{G}{x}
  &=\sum_{n=0}^{\infty}\left((1-x^2)P_n''(x)-2xP_n'(x)\right)t^n。
\end{align*}
右辺は：
\begin{align*}
  -t^2\nhb{2}{G}{t}-2t\hb{G}{t}
  &=-\sum_{n=0}^{\infty}n(n+1)P_n(x)t^n。
\end{align*}
$t^n$ の係数を比較する：
\begin{align*}
  (1-x^2)P_n''(x)-2xP_n'(x)+n(n+1)P_n(x)
  &=0。
\end{align*}
以上より、Legendre 多項式の満たす微分方程式が示された。

\subsection*{問題4}

問題の定義および各設問は2023年度問題4と同一である。したがって、同じ導出により：
\begin{align*}
  J_{-n}(z)
  &=(-1)^nJ_n(z),\\
  \bbn{J_n(z)}{z}
  &=\frac{1}{2}\left(J_{n-1}(z)-J_{n+1}(z)\right),\\
  \frac{n}{z}J_n(z)
  &=\frac{1}{2}\left(J_{n-1}(z)+J_{n+1}(z)\right),\\
  \nbbn{2}{J_n}{z}
  +\frac{1}{z}\bbn{J_n}{z}
  +\left(1-\frac{n^2}{z^2}\right)J_n
  &=0。
\end{align*}

\section*{2025年度}

\subsection*{問題1}

1次元 Poisson 方程式の Green 関数は：
\begin{align*}
  \nhb{2}{G(x,x')}{x}
  &=-\delta(x-x')
\end{align*}
を満たす。

\subsubsection*{1}

問題文の Fourier 表示：
\begin{align*}
  G(x,x')
  &=\frac{1}{2\pi}\int_{-\infty}^{\infty}g(k)e^{ik(x-x')}\,dk,\\
  \delta(x-x')
  &=\frac{1}{2\pi}\int_{-\infty}^{\infty}e^{ik(x-x')}\,dk
\end{align*}
を用いる。Green 関数を $x$ で2回微分する：
\begin{align*}
  \nhb{2}{G(x,x')}{x}
  &=-\frac{1}{2\pi}\int_{-\infty}^{\infty}k^2g(k)e^{ik(x-x')}\,dk。
\end{align*}
Poisson 方程式と比較すると：
\begin{align*}
  k^2g(k)&=1。
\end{align*}
したがって：
\begin{align*}
  g(k)&=\frac{1}{k^2}。
\end{align*}
ただし、$k=0$ の特異点があるため、逆変換には境界条件に応じた処方が必要である。

\subsubsection*{2}

$x<x'$ で零となる形として：
\begin{align*}
  G(x,x')
  &=A(x-x')\Theta(x-x')
\end{align*}
と置く。微分すると：
\begin{align*}
  \hb{G}{x}
  &=A\Theta(x-x'),\\
  \nhb{2}{G}{x}
  &=A\delta(x-x')。
\end{align*}
Poisson 方程式と比較して $A=-1$ を得る。したがって：
\begin{align*}
  G(x,x')
  &=-(x-x')\Theta(x-x') \\
  &=\begin{cases}
    0 & x<x',\\
    -(x-x') & x\geq x'
  \end{cases}。
\end{align*}

\subsection*{問題2}

区間 $0\leq x\leq L$, $0\leq t<\infty$ において、熱方程式：
\begin{align*}
  \hb{u(x,t)}{t}
  &=\kappa\nhb{2}{u(x,t)}{x}
\end{align*}
を、初期条件 $u(x,0)=f(x)$ および Neumann 境界条件：
\begin{align*}
  \hb{u(0,t)}{x}&=0,
  &\hb{u(L,t)}{x}&=0
\end{align*}
の下で解く。

\subsubsection*{1}

変数分離形 $u(x,t)=X(x)T(t)$ を代入する：
\begin{align*}
  X(x)T'(t)
  &=\kappa X''(x)T(t)。
\end{align*}
両辺を $\kappa X(x)T(t)$ で割る：
\begin{align*}
  \frac{1}{\kappa T(t)}\bbn{T(t)}{t}
  &=\frac{1}{X(x)}\nbbn{2}{X(x)}{x}
   =-\lambda。
\end{align*}
したがって：
\begin{align*}
  \nbbn{2}{X}{x}+\lambda X&=0,\\
  \bbn{T}{t}+\kappa\lambda T&=0。
\end{align*}

\subsubsection*{2}

$\lambda=k^2>0$ のとき：
\begin{align*}
  X(x)&=A\cos kx+B\sin kx。
\end{align*}
$X'(0)=0$ から $B=0$ を得る。さらに $X'(L)=0$ から：
\begin{align*}
  \sin kL&=0,
  &k_n&=\frac{n\pi}{L}
\end{align*}
を得る。したがって：
\begin{align*}
  X_n(x)&=\cos\frac{n\pi x}{L},\\
  T_n(t)&=\exp\left[-\kappa\left(\frac{n\pi}{L}\right)^2t\right]
  \qquad (n\geq1)。
\end{align*}
また、$\lambda=0$ のときは $X_0(x)=1$, $T_0(t)=1$ である。$\lambda<0$ では Neumann 境界条件を満たす非自明解は存在しない。

\subsubsection*{3}

各固有モードを重ね合わせる：
\begin{align*}
  u(x,t)
  &=A_0+\sum_{n=1}^{\infty}A_n
    \exp\left[-\kappa\left(\frac{n\pi}{L}\right)^2t\right]
    \cos\frac{n\pi x}{L}。
\end{align*}
$t=0$ で初期条件を課すと、$f(x)$ の Fourier 余弦級数を得る。直交関係から：
\begin{align*}
  A_0
  &=\frac{1}{L}\int_0^Lf(\xi)\,d\xi,\\
  A_n
  &=\frac{2}{L}\int_0^Lf(\xi)\cos\frac{n\pi\xi}{L}\,d\xi
  \qquad (n\geq1)。
\end{align*}
したがって：
\begin{align*}
  u(x,t)
  &=\frac{1}{L}\int_0^Lf(\xi)\,d\xi \\
  &\quad+\frac{2}{L}\sum_{n=1}^{\infty}
    \left(\int_0^Lf(\xi)\cos\frac{n\pi\xi}{L}\,d\xi\right)
    \exp\left[-\kappa\left(\frac{n\pi}{L}\right)^2t\right]
    \cos\frac{n\pi x}{L}。
\end{align*}

\subsubsection*{4}

$n\geq1$ のモードは $t\to+\infty$ で減衰するため：
\begin{align*}
  \lim_{t\to+\infty}u(x,t)
  &=\frac{1}{L}\int_0^Lf(\xi)\,d\xi。
\end{align*}
これは初期温度の空間平均であり、十分な時間の後には一様な温度分布へ近づく。実際、Neumann 境界条件により：
\begin{align*}
  \bbn{}{t}\int_0^Lu(x,t)\,dx
  &=\kappa\left[\hb{u}{x}\right]_0^L
   =0
\end{align*}
であるから、区間内の全熱量は保存される。

\subsection*{問題3}

Legendre 多項式を Rodrigues の公式：
\begin{align*}
  P_n(x)
  &=\frac{1}{2^nn!}\nbbn{n}{}{x}(x^2-1)^n
\end{align*}
によって定義する。

\subsubsection*{1}

Rodrigues の公式へ $n=0,1,2$ を代入する：
\begin{align*}
  P_0(x)&=1,\\
  P_1(x)&=x,\\
  P_2(x)&=\frac{1}{2}(3x^2-1)。
\end{align*}

\subsubsection*{2}

$0\leq m<n$ とする。Rodrigues の公式を代入する：
\begin{align*}
  \int_{-1}^{1}x^mP_n(x)\,dx
  &=\frac{1}{2^nn!}\int_{-1}^{1}x^m\nbbn{n}{}{x}(x^2-1)^n\,dx。
\end{align*}
部分積分を $n$ 回行う。$(x^2-1)^n$ は $x=\pm1$ に $n$ 重の零点を持つため、すべての境界項は消える：
\begin{align*}
  \int_{-1}^{1}x^mP_n(x)\,dx
  &=\frac{(-1)^n}{2^nn!}\int_{-1}^{1}\nbbn{n}{x^m}{x}(x^2-1)^n\,dx。
\end{align*}
$m<n$ では $\nbbn{n}{x^m}{x}=0$ である。したがって：
\begin{align*}
  \int_{-1}^{1}x^mP_n(x)\,dx
  &=0
  \qquad (0\leq m<n)。
\end{align*}

\subsubsection*{3}

まず、$P_n(x)$ の最高次係数を求める。Rodrigues の公式における最高次項から：
\begin{align*}
  P_n(x)
  &=\frac{(2n)!}{2^n(n!)^2}x^n+\text{低次の項}。
\end{align*}
最高次係数を：
\begin{align*}
  c_n&=\frac{(2n)!}{2^n(n!)^2}
\end{align*}
とする。$x^n=P_n(x)/c_n+q_{n-1}(x)$ と書け、問2の直交性から：
\begin{align*}
  \int_{-1}^{1}P_n(x)^2\,dx
  &=c_n\int_{-1}^{1}x^nP_n(x)\,dx。
\end{align*}
右辺の積分へ Rodrigues の公式を代入し、部分積分を $n$ 回行う：
\begin{align*}
  \int_{-1}^{1}x^nP_n(x)\,dx
  &=\frac{1}{2^n}\int_{-1}^{1}(1-x^2)^n\,dx。
\end{align*}
ここで、$x=2y-1$ と置く：
\begin{align*}
  \int_{-1}^{1}(1-x^2)^n\,dx
  &=2^{2n+1}\int_0^1y^n(1-y)^n\,dy \\
  &=2^{2n+1}\frac{(n!)^2}{(2n+1)!}。
\end{align*}
したがって：
\begin{align*}
  \int_{-1}^{1}P_n(x)^2\,dx
  &=\frac{(2n)!}{2^n(n!)^2}
    \frac{2^{n+1}(n!)^2}{(2n+1)!} \\
  &=\frac{2}{2n+1}。
\end{align*}

\subsection*{問題4}

Bessel 関数 $J_0(x)$ を：
\begin{align*}
  J_0(x)
  &=\frac{1}{2\pi}\int_0^{2\pi}e^{ix\cos\theta}\,d\theta
\end{align*}
で定義する。Laplace 変換を：
\begin{align*}
  L(s)
  &=\int_0^\infty e^{-sx}J_0(x)\,dx
\end{align*}
とする。ただし、$\operatorname{Re}s>0$ とする。

\subsubsection*{1}

固定した $\theta$ に対して：
\begin{align*}
  \int_0^\infty e^{-sx}e^{ix\cos\theta}\,dx
  &=\int_0^\infty e^{-(s-i\cos\theta)x}\,dx \\
  &=\frac{1}{s-i\cos\theta}。
\end{align*}

\subsubsection*{2}

積分順序を交換し、問1の結果を用いる：
\begin{align*}
  L(s)
  &=\frac{1}{2\pi}\int_0^{2\pi}\frac{d\theta}{s-i\cos\theta}。
\end{align*}
$z=e^{i\theta}$ と置く。このとき：
\begin{align*}
  d\theta&=\frac{dz}{iz},
  &\cos\theta&=\frac{1}{2}\left(z+\frac{1}{z}\right)。
\end{align*}
したがって：
\begin{align*}
  L(s)
  &=\frac{1}{2\pi}\oint_{|z|=1}\frac{2\,dz}{z^2+2isz+1}。
\end{align*}
分母の零点は：
\begin{align*}
  z_\pm
  &=-is\pm i\sqrt{s^2+1}。
\end{align*}
$s>0$ では：
\begin{align*}
  z_+
  &=i\left(\sqrt{s^2+1}-s\right)
\end{align*}
のみが単位円内にある。被積分関数の $z=z_+$ における留数は：
\begin{align*}
  \mathop{\mathrm{Res}}_{z=z_+}\frac{2}{(z-z_+)(z-z_-)}
  &=\frac{2}{z_+-z_-}
   =\frac{1}{i\sqrt{s^2+1}}。
\end{align*}
留数定理を用いる：
\begin{align*}
  L(s)
  &=\frac{1}{2\pi}\,2\pi i\,
    \frac{1}{i\sqrt{s^2+1}} \\
  &=\frac{1}{\sqrt{s^2+1}}。
\end{align*}
以上より：
\begin{align*}
  \int_0^\infty e^{-sx}J_0(x)\,dx
  &=\frac{1}{\sqrt{s^2+1}}
  \qquad (\operatorname{Re}s>0)
\end{align*}
である。

\end{document}
